\subsection{Arbitrary twists, constant local evaluations, and the global obstruction}

\begin{statement}{A.128}{Proved}
{All local common-character evaluations on the two actual sources}
Let $c_1,\ldots,c_k$ be distinct rational numbers, let
$a_1,\ldots,a_k\in\Q^\times$, and let $4\mid N$.
Use the actual old cover, new boundary fibre, and common field
\[
 \begin{gathered}
 C_N:\quad T-c_i=a_iX_i^N,\qquad X_i\ne0,\\
 c\in\Q\setminus\{c_1,\ldots,c_k\},\qquad
 D_i=(c-c_i)/a_i,\\
 \mathcal E_{N,c}
   =\Q[Z_i^{\pm1}:1\le i\le k]/(Z_i^N-D_i)_i,\\
 B_N=\text{the largest abelian field in every factor of }\mathcal E_{N,c},\\
 K=\Q(\sqrt{D_i}:1\le i\le k).
 \end{gathered}
\tag{A.128.1}
\]
For every common character $\chi$ of
$\operatorname{Gal}(B_N/\Q)$, the evaluation of
$(T-c,\chi)$ is constant on the ordinary points
$C_N(\Q_p)$ with $T\ne c$, for every odd prime $p$,
and on the corresponding real points.
There is no restriction on the number or reduction of the old punctures.

When the old finite punctures are either of the two actual sets
\[
       \{0,1,\tfrac12\}
             \qquad\text{or}\qquad
       \{0,1,\tfrac23,-\tfrac13\},
\tag{A.128.2}
\]
the same constancy holds at $p=2$.
The coefficients $a_i$ are arbitrary rational numbers;
the assertion does not require a rational point on $C_N$.

For either set in (A.128.2), let $\overline C_N$ be the
smooth proper compactification.
Each cyclic class extends to
\[
              \alpha_\chi=(T-c,\chi)\in\operatorname{Br}(\overline C_N).
\tag{A.128.3}
\]
Its evaluation is constant on all of $\overline C_N(\Q_v)$.
For an everywhere locally soluble old cover, the finite sum
\[
 \begin{gathered}
 \beta_{C_N,c}(\chi)=
     \sum_v\operatorname{inv}_v\alpha_\chi(P_v),\\
 \chi(\gamma_{C_N,c})=\beta_{C_N,c}(\chi),\qquad
            \gamma_{C_N,c}\in\operatorname{Gal}(B_N/K)
 \end{gathered}
\tag{A.128.4}
\]
is independent of every choice of local points $P_v$.
It is the actual Brauer--Manin evaluation of the homomorphism
$\chi\mapsto\alpha_\chi$.
The following are equivalent:
\begin{enumerate}
\item $\gamma_{C_N,c}=1$.
\item Some $a_c\in\Q^\times$ makes the augmented cover
      $T-c=a_cY^N$, $Y\ne0$, everywhere locally soluble.
\item Every finite collection consisting of one ordinary old point
      at each selected place, with $T\ne c$, lifts to an
      everywhere locally soluble such augmented cover.
\end{enumerate}
A rational point of the proper old curve makes this obstruction zero.
For arbitrary twists the formula retains the global obstruction:
the local constancy assertion alone does not evaluate (A.128.4).

The obstruction is invariant under the explicit coefficient
isomorphisms and changes of coordinates below.
For compatible old coefficient classes at exponents $N\mid M$,
the actual power maps give
$B_N\subset B_M$ and
$\gamma_{C_M,c}|_{B_N}=\gamma_{C_N,c}$.
Consequently an ordinary completed old tuple has a compatible
family of these finite obstruction elements, and the existence
of a completed new coordinate forces every one of them to be trivial.
\end{statement}
\begin{proof}
\emph{The two-point map and every relation character.}
Fix a local field $k_v=\Q_v$ and two ordinary old points
$P_R=(R,(\xi_i)_i)$ and $P_T=(T,(\eta_i)_i)$, with
$R,T\ne c$.
There is nothing to prove when $T=R$.
Otherwise put
\[
 x=\frac{T-c}{R-c},\qquad w_i=\frac{\eta_i}{\xi_i},
       \qquad d_i=\frac{D_i}{\xi_i^N}=\frac{c-c_i}{R-c_i}.
 \quad\text{Then}\quad
                     w_i^N=x+(1-x)d_i.
\tag{A.128.5}
\]
Thus constancy is exactly $(x,\chi)=0$.
Write $N=nv$, where $n=2^r\ge4$ and $v$ is odd.
The morphism $C_N\to C_n$ sends $X_i$ to $X_i^v$,
leaves $T$ and every $D_i$ unchanged, and induces
$Z_i\mapsto Z_i^v$ on the boundary roots.
A.126.21 proves $B_N=B_n$ for this actual map.
We can therefore work at exponent $n$.

At a growing step put $m=n/2$.
Every relation and its exact local image are
\[
 \begin{gathered}
 \prod_iD_i^{t_i}=\sigma r_t^m,\qquad
   0\le t_i<m,\quad \sigma\in\{1,-1\},\quad r_t\in\Q_{>0},\\
 \prod_id_i^{t_i}=\sigma\rho^m,\qquad
             \rho=\frac{r_t}{\prod_i\xi_i^{2t_i}}.
 \end{gathered}
\tag{A.128.6}
\]
The displayed correction is a local square.
Hence the quadratic character of $\rho$ is precisely the
restriction of the character of $r_t$.
For the original exponent $N$, replace $\xi_i$ here by
$\xi_i^v$; the correction remains a square.
The complete A.126 recursion says that, when
$\mu_n\subset B_m$, the new characters are
\[
        \chi_\rho\quad(\sigma=1),\qquad
        \chi_{2n}+\chi_\rho\quad(\sigma=-1),
       \qquad \chi_{2n}(5)=1/m,\quad\chi_{2n}(-1)=0,
\tag{A.128.7}
\]
together with the lower characters.
These are the restrictions of the original global characters,
with their literal correction (A.128.6).
In particular the local numbers $d_i$ are not substituted
for the rational input of the global common-field theorem.

The exponent bound furnished by its actual ambient field is
\[
          \operatorname{exp}\operatorname{Gal}(B_n/\Q)\mid n/2,
       \qquad B_n\subset\Q(\mu_{2n},\sqrt P:P\in\mathcal P),
       \quad
       \mathcal P=\{P:v_P(D_i)\ne0\text{ for some }i\}.
\tag{A.128.8}
\]
At exponent four, A.118 proves all the asserted odd-place
evaluations for arbitrary coefficients.
At each later step the coordinate-square map
$C_n\to C_m$ retains $x$.
Thus induction gives zero evaluation for every character
of $B_m$.
If $\mu_n\not\subset B_m$, the stopping rule A.126 gives
$B_n=B_m$, and there is nothing further to test.
At a growing step we have
\[
              \mu_n\subset B_m,\qquad
              i\in K,\qquad
              k_v(\sqrt{d_i}:i)=k_vK.
\tag{A.128.9}
\]
For the middle assertion, $i\notin K$ would make the
entire tower equal to $K$ by A.126's first stopping rule.
The last equality uses $D_i/d_i=\xi_i^n$, an actual square.
All cyclotomic characters of $\Q(\mu_n)$ therefore have
zero evaluation on $x$.

\emph{All odd primes, with arbitrary old punctures.}
We now prove that every character (A.128.7) evaluates
to zero over $k=\Q_p$, with $p$ odd.
We use the local reciprocity norm isomorphism in
\cite[Theorem 4.1.2]{BIR_KedlayaCFT}.
The characters of $2$-power cyclotomic fields are
unramified at $p$; their evaluation on $x$ is their
Frobenius value multiplied by $h=v_p(x)$.
The lower cyclotomic vanishing gives
\[
                         \operatorname{ord}_{n}(p)\mid h.
\tag{A.128.10}
\]
The following exact inversion permits $h\ge0$:
\[
 x'=x^{-1},\quad w_i'=w_i^{-1},\quad
 d_i'=d_i/w_i^n,\quad
                    \rho'=\rho/\prod_iw_i^{2t_i}.
\tag{A.128.11}
\]
It preserves (A.128.5)--(A.128.6), the local squareclass
field in (A.128.9), and every second character.
It negates its cyclic evaluation, which preserves vanishing.

Suppose first that $p\equiv1\pmod n$.
Then $\mu_n\subset k$.
For a nonsquare $x$ and $\theta^n=x$, the field
$L=k(\theta)$ is a splitting field whose Galois group
embeds in $\mu_n$.
A proper subgroup would fix $\theta^{n/2}=\sqrt x$.
Thus $L/k$ is cyclic of degree $n$, with quadratic
subfield $k(\sqrt x)$, and
\[
       d_i=N_{L/k}\left(\frac{w_i-\theta}{1-\theta}\right).
\tag{A.128.12}
\]
In a positive relation, $m$ times the reciprocity image
of $\rho$ is zero in $\Z/n$.
This makes that image even, which is precisely the kernel
of the restriction to the quadratic subfield.
Therefore $(x,\rho)=0$.
In a negative relation use instead
$(\zeta_n\rho)^m=-\rho^m$ in the same norm calculation.
The restriction of $\chi_{2n}$ to $G_k$ is exactly
the quadratic character of $\zeta_n$: both cut out
$k(\mu_{2n})/k$, with value $1/2$ on its nontrivial
automorphism when that extension is nontrivial.
Hence $(x,\chi_{2n}+\chi_\rho)=0$ as well.
For square $x$ all quadratic evaluations are zero;
$h$ is even and this restricted cyclotomic character
has order at most two, so its evaluation is also zero.

For the remaining cases write $x=p^hu$ with $u$ a unit,
and put $s_i=nv_p(w_i)$.
When $h\notin n\Z$ one has $h>0$, none of the $s_i$
equals $h$, and $1-x$ is a unit with residue one.
The valuation and normalized unit of $d_i$ are respectively
$(s_i,\text{an }n\text{th-power residue})$ if $s_i<h$,
and $(h,-\bar u)$ if $s_i>h$.
With
$b=\sum_{s_i>h}t_i$, the relation in (A.128.6) gives
\[
 m v_p(\rho)=\sum_{s_i<h}t_is_i+hb,
                 \qquad m v_p(\rho)\equiv hb\pmod n.
\tag{A.128.13}
\]

Take $p\equiv1\pmod4$, and set $t=v_2(p-1)<r$ and
$d=2^{r-t}$.
Successive squaring gives
$v_2(p^{2^j}-1)=t+j$.
Thus (A.128.10) says $h=dH$.
Also $v_2(5^j-1)=2+v_2(j)$ shows that
$\chi_{2n}(p)$ has order $2d$ and odd numerator with
that denominator.  Consequently
\[
                       h\chi_{2n}(p)=H/2\quad\text{in }\Q/\Z.
\tag{A.128.14}
\]
Here $h$ is even.

For nonsquare $x$, the residue $\bar u$ is a nonsquare.
If $h\in n\Z$, the leading residues of $w_i^n$ and $x$
cannot cancel when their valuations agree.
The denominator $1-x$ is a unit, also when $h=0$,
because then $\bar u\ne1$.
It follows that every $v_p(d_i)$ is in $n\Z$, and
$v_p(\rho)$ is even.
The quadratic extension $k(\sqrt x)/k$ is unramified,
so $\rho$ is a norm; the cyclotomic evaluation is zero
since $h$ is a multiple of $n$.

For nonsquare $x$ with $h\notin n\Z$, choose a residue
character $\eta:\F_p^\times\to\mu_{2^t}$ of full order.
One has $\eta(-1)=-1$.
Since $t\ge2$ and $\bar u$ is a nonsquare,
$\eta(-\bar u)$ is a primitive $2^t$th root.
The normalized-unit relation, whose right-hand
$m$th power is killed by $\eta$, gives
\[
 b\equiv
 \begin{cases}
 0\pmod{2^t},&\sigma=1,\\
 2^{t-1}\pmod{2^t},&\sigma=-1.
 \end{cases}
\tag{A.128.15}
\]
Equation (A.128.13) now makes $v_p(\rho)$ even in
the positive case, and congruent to $H$ modulo two
in the negative case.
Since $k(\sqrt x)$ is unramified quadratic, its
evaluation on $\rho$ is respectively $0$ and $H/2$.
The latter cancels exactly (A.128.14).

For square $x$ all positive characters vanish.
A negative character is also zero when $H$ is even.
If $H$ were odd, $h\notin n\Z$ and (A.128.13)
would force $2^{t-1}\mid b$.
But $\bar u$ is a square, so $\eta(-\bar u)$ has order
dividing $2^{t-1}$.
Its $b$th power is then $1$, whereas the negative
normalized-unit relation requires $\eta(-1)=-1$.
No such relation exists in this case.
This proves every evaluation for $p\equiv1\pmod4$.

Finally take $p\equiv3\pmod4$.
The lower character of $i$ gives even $h$.
For nonsquare $x$ and $h\notin n\Z$, $\bar u$ is
nonsquare, so $-\bar u$ is square.
The valuation alternatives above show that every $d_i$
has even valuation and square normalized residue.
Every principal unit is a square by Hensel lifting,
so all $d_i$ are squares in $k$.
This contradicts the literal field inclusion
$k(i)\subset kK=k(\sqrt{d_i}:i)$ in (A.128.9).
Thus $h\in n\Z$.
The same noncancellation argument as above gives
$v_p(\rho)$ even and both evaluations zero.

For square $x$, only a negative relation needs checking.
If $h\notin n\Z$, the quadratic residue character of
the normalized-unit relation makes $b$ odd,
because $-\bar u$ is nonsquare.
Equation (A.128.13) then forces $m\mid h$.
As $\chi_{2n}$ has order dividing $m$, its evaluation
is zero.  For $h\in n\Z$ it is zero directly.
This exhausts all odd primes and proves the induction
there for every finite set of old punctures.

At the real place all old ratios in (A.128.5) are
positive, so $R,T$ lie in the same interval relative
to every old puncture.
When $c$ is outside the interval between $R$ and $T$,
$x>0$ and every real cyclic evaluation is zero.
When $c$ is between them, every $d_i>0$ and hence
$D_i=d_i\xi_i^n>0$.
The original boundary algebra then has a real root tuple.
Every common character restricts trivially to its
stabilizer $G_\R$, proving the real assertion.

\emph{Exact affine and inverse-coordinate maps for arbitrary twists.}
An affine map $g(z)=az+b$, $a\ne0$, sends
$c_i$ to $g(c_i)$, $a_i$ to $aa_i$, and leaves
the root coordinates unchanged.
It preserves each boundary radicand and $x$, while
\[
                      (g(T)-g(c),\chi)=(T-c,\chi)+(a,\chi).
\tag{A.128.16}
\]
For an old puncture $c_0$ and $b\ne0$, the inverse-coordinate
map $U=b/(T-c_0)$ has the following actual root equations:
\[
 \begin{gathered}
 a'_0=b/a_0,\quad Y_0=X_0^{-1},\\
 c'_i=b/(c_i-c_0),\quad
 a'_i=-\frac{ba_i}{a_0(c_i-c_0)},\quad
 Y_i=X_i/X_0\quad(i\ne0),\\
 U=a'_0Y_0^n,\qquad U-c'_i=a'_iY_i^n.
 \end{gathered}
\tag{A.128.17}
\]
The inverse is $T=c_0+b/U$, $X_0=Y_0^{-1}$,
$X_i=Y_i/Y_0$.
At the new boundary it induces the algebra isomorphism
\[
 \begin{gathered}
 Z'_0=Z_0^{-1},\quad Z'_i=Z_i/Z_0,\qquad
 D'_0=D_0^{-1},\quad D'_i=D_i/D_0,\\
 x'=x\,\frac{R-c_0}{T-c_0}
        =x(\xi_0/\eta_0)^n,\\
 (U-g(c),\chi)=(T-c,\chi)+
                     \left(-\frac b{a_0(c-c_0)},\chi\right).
 \end{gathered}
\tag{A.128.18}
\]
These maps identify every actual field factor and common
character, commute with coordinate-power maps, and preserve
the normalized evaluation by (A.128.8).
In particular they transport every mixed character, not
only the characters of individual coordinates.

\emph{The dyadic base at exponent four.}
For the harmonic set, the individual, singleton, and
two-coordinate common characters of A.96 have constant
classes already on the corresponding one- or two-coordinate
cover.  The forgetful maps and (A.96.18) pull those
identities back to the present cover.
For a negative triple relation, (A.128.6) gives the
same local character with $\rho=r_t/\prod_i\xi_i^2$.
The two-point ratios and this exact local relation satisfy
the hypotheses of A.111, so its normalized evaluation
is zero.
For the cubic set, A.120 applies to every negative subset
after exactly the same square correction.
Together with the complete quartic common-field formula
A.96, these prove the dyadic base for arbitrary coefficients.

\emph{The dyadic induction and its actual coordinate reductions.}
At a later growing step retain (A.128.9).
In particular a subset product of the rational $D_i$
is a negative rational square.
The corresponding local subset product of the $d_i$
is still a negative square by (A.128.5).
The unit power and signed cyclotomic norm identities are
\[
 (\Z_2^\times)^{2^j}=1+2^{j+2}\Z_2\quad(j\ge1),\qquad
 \operatorname{rec}_2(u)|_{\mu_{2^r}}:
                   \zeta\longmapsto\zeta^{u^{-1}}.
\tag{A.128.19}
\]
The first follows by iterating the polynomial squaring
map (A.127.7); the second is the actual local reciprocity
map of \cite[Example 4.1.4]{BIR_KedlayaCFT}.
Thus lower vanishing on a unit $x$ implies
$x\in1+2^r\Z_2$, with the sign retained.

If the valuation coordinates of $R,T$ at all old punctures
agree, put $M=\max_i v_2(R-c_i)$ and $s=v_2(T-R)$.
The old $n$th-power unit ratios give $s\ge M+r+2$.
If $v_2(c-R)\ge M+2$, every $d_i$ lies in $1+4\Z_2$,
contrary to the negative-square subset.
It follows that
\[
                    x\in1+2^{r+1}\Z_2
                              =(\Z_2^\times)^{n/2}.
\tag{A.128.20}
\]
This kills all the characters in (A.128.8).
This part of the argument holds for any old puncture set.

For the harmonic set, apply the maps $1-T$ and $1/(2T)$
using (A.128.16)--(18).
They put $v_2(R)=e\ge1$ and preserve the source punctures.
The exact ratio argument (A.127.12) puts $v_2(T)=e_T>0$.
For unequal depths it gives
$e,e_T\ge r+2$ and $e\equiv e_T\pmod n$.
The other finite punctures have nonpositive valuation.

For the cubic set, every $n$th power is a fourth power,
so the exhaustive valuation table (A.120.11) and its
two cross-row contradictions (A.120.12)--(13) apply
without a coefficient restriction: they concern only
the literal ratios $(T-c_i)/(R-c_i)$.
Thus unequal valuation coordinates occur in one boundary
row.  For a finite row near $c_i$, let
$\delta=\max_{j\ne i}v_2(c_i-c_j)$, which is $1$
at $0,2/3$ and $2$ at $1,-1/3$.
The ratio at a nearest other puncture is an
$n$th-power unit.  At unequal depths its distance
from $1$ has valuation
$\min(v_2(R-c_i),v_2(T-c_i))-\delta$.
Equation (A.128.19) therefore gives
\[
             \min(v_2(R-c_i),v_2(T-c_i))\ge\delta+r+2.
\tag{A.128.21}
\]
The actual affine map $(T-c_i)/2^\delta$ gives old
boundary zero, all other boundaries of nonpositive
valuation, and both depths at least $r+2$.
For the infinite row, the unit ratio
$(1-1/T)/(1-1/R)$ gives
$\min(-v_2(R),-v_2(T))\ge r+2$ at unequal depths.
The actual map (A.128.17) with $c_0=0,b=1$ has
new finite punctures $0,1,3/2,-3$ and gives the
same depth conclusion.

We have thus reached the following general configuration:
old boundaries zero and numbers $b$ with $v_2(b)\le0$,
and
\[
 \begin{gathered}
 e=v_2(R),\quad e_T=v_2(T),\quad
 e,e_T\ge r+2,\quad e\ne e_T,\\
 e\equiv e_T\pmod n,\quad s=v_2(T-R)=\min(e,e_T).
 \end{gathered}
\tag{A.128.22}
\]
All coefficients, common fields, lower characters, and
evaluations have been carried by the explicit maps above.

Put $k=v_2(c)$.
For $k\le0$, $v_2(x-1)=s-k\ge r+2$, so (A.128.20)
applies.
For $1\le k\le r+1$, $x$ is a unit and
$v_2(x-1)=s-k$.
Lower cyclotomic vanishing and (A.128.19) force
$s-k\ge r$.
A depth at least $r+1$ again gives (A.128.20).
The only remaining case is
\[
          s-k=r,\qquad k\ge2,\qquad
          e-k\equiv r\pmod n,\qquad
          x\in(\Z_2^\times)^{n/4}.
\tag{A.128.23}
\]
Here $v_2(d_0)=k-e$, and every other
$d_b=(1-c/b)/(1-R/b)$ lies in $1+4\Z_2$.
Each original global relation has exactly the local
image (A.128.6).
Taking its valuation gives $m\mid t_0(k-e)$.
Since $v_2(k-e)=v_2(r)<r-1$, $t_0$ is even.
After removing the equal powers of two, the left
unit in that relation is $1$ modulo four.
A negative right side would be $-1$ modulo four,
because $m$ is even.
Thus every original relation in the complete A.126
list has positive sign.
That formula makes $B_n$ the compositum of $B_m$
and quadratic fields.
Its exponent consequently divides $n/4$, and the
last power assertion in (A.128.23) kills all its
characters.  This uses the global sign in (A.128.6);
no new common field for nonrational radicands is invoked.

Finally let $k\ge r+2$.
For each nonzero old boundary $b$,
$d_b\in1+2^{r+2}\Z_2$ is an actual $n$th power,
say $d_b=u_b^n$.
The original radicand is
$D_b=(\xi_bu_b)^n$.
Specializing the original boundary coordinates gives
the actual local algebra surjection
\[
 \mathcal E_{n,c}\otimes\Q_2
       \longrightarrow \Q_2[Z_0]/(Z_0^n-D_0),
                          \qquad Z_b\longmapsto\xi_bu_b.
\tag{A.128.24}
\]
Each factor $E$ of the target contains the local cyclic
field of every common character.
Indeed its root tuple defines a local stabilizer contained
in a global root stabilizer, on which that character
vanishes.
For one root $\theta_E$ in each factor the literal
polynomial norm identity is
\[
 x=\frac{\eta_0^n-D_0}{\xi_0^n-D_0}
      =\prod_E N_{E/\Q_2}
              \left(\frac{\eta_0-\theta_E}{\xi_0-\theta_E}\right).
\tag{A.128.25}
\]
All factors are defined since $R,T\ne c$.
Norm transitivity makes this a norm from every character
field.  This finishes the dyadic induction for both
sources with arbitrary rational coefficients.
The odd-factor map after (A.128.5) proves the assertion
for every $4\mid N$.

\emph{The actual proper Brauer class and the global constant sum.}
The old cover is geometrically integral.
For completeness, over $\overline\Q(T)$ the radicands
$(T-c_i)/a_i$ are independent modulo $N$th powers:
the valuation at $T=c_i$ detects its exponent modulo $N$.
Equivalently, when adjoining them successively, the
previous roots are unramified at the next $c_i$, where
its new radicand has valuation one; its root therefore
has degree $N$ by the valuation criterion.
Thus the geometric degree is $N^k$.
The ordinary curve is smooth, since each nonzero-root
equation has invertible derivative in characteristic zero.

Compute the residues of $\alpha_\chi$ on its smooth
proper model.
At $T=c$ the cover is etale and its closed residue
fields are precisely the factors of $\mathcal E_{N,c}$.
The character restricts to zero in every such field,
so the residue $1\cdot\chi$ is zero.
At each old finite puncture, $T-c$ is a unit.
At infinity put $z=1/T$ and pass to a finite constant
extension containing $\mu_N$ and the $N$th roots of
the $a_i$.
In formal Laurent series the root equations become
\[
        X_i=a_i^{-1/N}z^{-1/N}(1-c_i z)^{1/N}.
\tag{A.128.26}
\]
The final factors are actual formal unit roots.
Thus the geometric ramification index is exactly $N$,
and the constant extension is unramified for this
valuation.  Consequently $v_P(T-c)=-N$ at every
point above infinity.
Its residue is $-N\chi=0$, by (A.128.8).
All remaining residues are zero because $T-c$ is a unit.
The residue formula and unramified exact sequence
\cite[\S1.4 and Corollary 3.7.3]{AgentCTSkoro2019}
prove (A.128.3).

The ordinary open set avoiding $T=c$ is dense in
every nonempty set of local points of this smooth
curve: an etale local parameter gives a local
one-dimensional ball, from which only finitely
many points have been removed.
Brauer evaluation is locally constant, including
over $\R$, by
\cite[Proposition 9.5.2]{AgentCTSkoro2019}.
The proved ordinary-point constancy therefore extends
to every point of $\overline C_N(\Q_v)$.
In particular proper local solubility supplies ordinary
local points whenever they are needed in A.85.

For a proper variety, a fixed Brauer class has zero
evaluation at all local points outside a finite set
of places
\cite[Proposition 12.3.1(iii)]{AgentCTSkoro2019}.
There are only finitely many common characters.
This proves that (A.128.4) is a finite sum uniformly
in the chosen points.
Linearity of the cyclic class makes it a homomorphism
on the finite character group.
Perfect duality for a finite abelian group supplies
the unique element $\gamma_{C_N,c}\in
\operatorname{Gal}(B_N/\Q)$ with the stated values.
This is precisely the Brauer--Manin pairing pulled
back by the actual map $\chi\mapsto\alpha_\chi$;
global constant classes have zero sum by
\cite[Theorem 12.1.8]{AgentCTSkoro2019}.

For each individual character the exact norm identity
\[
 \begin{gathered}
 T-c=a_i(X_i^N-D_i),\\
 X_i^N-D_i
    =N_{\Q(\sqrt{D_i})(C_N)/\Q(C_N)}
                     (X_i^{N/2}-\sqrt{D_i}),\qquad
                (T-c,D_i)=(a_i,D_i)
 \end{gathered}
\tag{A.128.27}
\]
holds, also in the split quadratic etale algebra.
Thus $\beta_{C_N,c}$ vanishes on every character of $K$.
Its dual element lies in the exact kernel
$\operatorname{Gal}(B_N/K)$, as asserted.

The finite reciprocity element of A.85 equals
$\gamma_{C_N,c}$: its evaluations on every character
are exactly (A.128.4), and its outside-set evaluations
are zero by that theorem's explicit calculation.
Local constancy lets its controlled old points be
chosen arbitrarily.
The necessity and sufficiency directions of A.85
therefore prove all three equivalences in the statement.
For a rational point $P\in\overline C_N(\Q)$, the
proper Brauer class has the global evaluation
$\alpha_\chi(P)\in\operatorname{Br}(\Q)$.
Its image at every completion is the same as the
evaluation at every local old point, by the proved
local constancy and functoriality of evaluation.
Global reciprocity makes their sum zero.
This includes rational boundary points and proves
the full rational-point assertion.
More generally $B_N=K$ forces zero by (A.128.27).
For the remaining twists (A.128.4) is the retained,
fully specified global obstruction.
Only the constancy of evaluations has been proved
locally; no local constant-class or Tate--Shafarevich
identification is used here.

\emph{Coefficient invariance and compatible completed maps.}
Replacing $a_i$ by $a_i u_i^N$ gives the actual
isomorphism $X_i\mapsto X_i/u_i$, and the same
map on boundary roots.
It leaves $T-c$ unchanged, so it preserves $B_N$,
the classes, and $\gamma_{C_N,c}$.
Equations (A.128.16) and (A.128.18) show that the
affine and inverse-coordinate maps change every
cyclic class only by its displayed global constant.
Its reciprocity sum is zero.
They too preserve the obstruction under the actual
boundary-field identifications.

Now suppose $N\mid M$, both divisible by four,
and representatives of compatible old classes satisfy
$a_i(M)=a_i(N)u_i^N$ with $u_i\in\Q^\times$.
The maps are literally
\[
 \begin{gathered}
 C_M\longrightarrow C_N,\qquad
          T\longmapsto T,\quad X_i(N)=u_iX_i(M)^{M/N},\\
 \mathcal E_{N,c}\longrightarrow\mathcal E_{M,c},
                         \qquad Z_i(N)\longmapsto u_iZ_i(M)^{M/N}.
 \end{gathered}
\tag{A.128.28}
\]
Over $\overline\Q$ the boundary-root map is onto.
Each upper factor therefore contains the selected
lower factor, proving $B_N\subset B_M$.
The individual quadratic fields agree, since
$a_i(M)/a_i(N)$ is an even power.
Every lower class pulls back exactly to the same
expression $(T-c,\chi)$ on the upper cover.
Evaluation at an upper local point followed by its
image, together with the proved independence of
local choices, gives
\[
       \beta_{C_M,c}|_{\widehat{\operatorname{Gal}(B_N/\Q)}}
                  =\beta_{C_N,c},\qquad
                 \gamma_{C_M,c}|_{B_N}=\gamma_{C_N,c}.
\tag{A.128.29}
\]
Different choices of $u_i$ in (A.128.28) give deck
changes fixing $T$ and the same character evaluations.
Hence the restriction equality is independent of those
choices.

An ordinary completed old tuple supplies such compatible
coefficient classes and old points at every finite
exponent.
Even when its definition specifies only finite places,
these old curves have real points as well.
Order the real punctures increasingly.
The only impossible sign pattern has $c_i<c_j$ and
$a_i<0<a_j$ for some pair.
But at every finite place the old equations give the
literal norm equation
\[
 U^2-\frac{a_j}{a_i}V^2=\frac{c_j-c_i}{a_i},
       \qquad U=X_i^{N/2},\quad V=X_j^{N/2}.
\tag{A.128.30}
\]
The constant quaternion
$(a_j/a_i,(c_j-c_i)/a_i)$ would therefore split
at every finite place and be nonsplit at the real
place, where both entries are negative.
This contradicts global reciprocity.
Consequently the signs select one nonempty real
interval; for a $T$ in that interval avoiding $c$,
all $(T-c_i)/a_i$ are positive and have actual
real $N$th roots.
The compatible coefficient classes have the same
signs at every level.
Thus (A.128.29) defines the asserted compatible family
of finite obstruction elements.
A completed new coordinate gives, at each exponent,
an augmented rational coefficient with all its actual
local lifts.  A.85 then forces $\gamma_{C_N,c}=1$
for each $N$.
The argument supplies these necessary compatible
obstruction elements.  A.132 proves that all-place
finite matching for the same fixed local centres
already gives unique, automatically compatible new
coefficients.  The vanishing obtained here supplies
the finite prescribed-point assertion of A.85, whose
points outside the controlled set may vary.
A.133 constructs a harmonic completed old tuple
for which every such finite assertion holds, while
all-place matching fails already modulo two.
The restriction maps (A.132.9) specify this precise
failure of passage from finite prescriptions to the
fixed family at all places.
\end{proof}

The retained certificate
\path{code/arbitrary_twist_constancy_certificate.py}
checks the literal coordinate and boundary-algebra maps,
finite odd-prime mixed-relation cancellations at several
dyadic exponents, the dyadic depth implications, and the
coefficient power-map identities.
Its finite checks supplement the unrestricted induction above.
